% Include this source from the TeDiServe paper as the Artifact Evaluation appendix.
\section{Artifact Evaluation}
\label{sec:artifact-evaluation}

The TeDiServe artifact is a research fork of vLLM containing the implementation
of the scheduler, step-cost predictor, resource manager, and reconfiguration
logic described in \S\ref{sec:tediserve}. It targets the \emph{Artifacts
Functional} and \emph{Artifacts Available} badges. The archival release and
full instructions are available at \url{ARTIFACT-DOI-OR-URL}.

\paragraph{Evaluator workflow.}
The supported portable workflow requires Linux, Python~3.12, and no GPU. After
building the included vLLM fork and installing LightGBM, evaluators run
\texttt{bash artifact/run\_demo.sh --mode smoke --instances 2 --requests 8}.
It completes in under one minute on the reference CPU login node and reports a
machine-readable result showing that all requests complete and that both
simulated model instances receive work. This run executes TeDiServe's main
deadline-aware scheduler with a simulated executor; it does not load model
weights or run CUDA kernels. The scheduler's control code runs, but fixed
simulated tokens/confidence values mean that this mode does not validate model
quality or meaningful confidence-threshold adaptation. It is consequently a
functional demonstration of multi-instance execution, not a performance or
accuracy reproduction.

Evaluators with two or more NVIDIA GPUs can additionally run
\texttt{bash artifact/run\_demo.sh --mode gpu --instances 2 --requests 8}.
This starts TeDiServe with two TP=1 instances of LLaDA-8B-Instruct and sends
concurrent requests. The GPU MWE uses a GH200-compatible environment and
downloads the model on first use. Dynamic reconfiguration is disabled in both
MWE modes; Gurobi is required only for reconfiguration experiments.

\paragraph{Relationship to the evaluation.}
The artifact MWE demonstrates the central implementation path---deadline-aware
request scheduling across model instances. It intentionally does not rerun all
paper plots: those experiments require long trace replays, multi-GPU clusters,
and, where reconfiguration is enabled, a licensed Gurobi installation. The
paper's Section~5.2 and~5.4 measurements used GH200 GPUs, while the accuracy
measurements associated with Figure~9 used H100 GPUs; Figure~9 is retained as
research reference rather than a promised evaluator target. The artifact guide
maps each paper component to its source location and precisely states the
resources, commands, expected output, and unsupported workflows.
